\documentclass[10pt, letterpaper]{article}

% ------------------------------------------------------------------------
% Cyrillic and font support for pdfLaTeX
% ------------------------------------------------------------------------
\usepackage[T2A]{fontenc}
\usepackage[utf8]{inputenc}
\usepackage[russian,english]{babel}
\usepackage{mathptmx} % Times New Roman-like serif font with Cyrillic support

% Packages:
\usepackage[
    ignoreheadfoot,
    top=2 cm,
    bottom=2 cm,
    left=2 cm,
    right=2 cm,
    footskip=1.0 cm
]{geometry}
\usepackage{titlesec}
\usepackage{tabularx}
\usepackage{array}
\usepackage[dvipsnames]{xcolor}
\definecolor{primaryColor}{RGB}{0, 0, 0}
\usepackage{enumitem}
\usepackage{fontawesome5}
\usepackage{amsmath}
\usepackage[
    pdftitle={Резюме Шахриёра Джабборова},
    pdfauthor={Shaxriyor Jabborov},
    pdfcreator={LaTeX},
    colorlinks=true,
    urlcolor=primaryColor
]{hyperref}
\usepackage[pscoord]{eso-pic}
\usepackage{calc}
\usepackage{bookmark}
\usepackage{lastpage}
\usepackage{changepage}
\usepackage{paracol}
\usepackage{ifthen}
\usepackage{needspace}

% Some settings:
\raggedright
\AtBeginEnvironment{adjustwidth}{\partopsep0pt}
\pagestyle{empty}
\setcounter{secnumdepth}{0}
\setlength{\parindent}{0pt}
\setlength{\topskip}{0pt}
\setlength{\columnsep}{0.15cm}
\pagenumbering{gobble}

\titleformat{\section}{\needspace{4\baselineskip}\bfseries\large}{}{0pt}{}[\vspace{1pt}\titlerule]

\titlespacing{\section}{
    -1pt
}{
    0.3 cm
}{
    0.2 cm
}

\renewcommand\labelitemi{$\vcenter{\hbox{\small$\bullet$}}$}

\newenvironment{highlights}{
    \begin{itemize}[
        topsep=0.10 cm,
        parsep=0.10 cm,
        partopsep=0pt,
        itemsep=0pt,
        leftmargin=0 cm + 10pt
    ]
}{
    \end{itemize}
}

\newenvironment{highlightsforbulletentries}{
    \begin{itemize}[
        topsep=0.10 cm,
        parsep=0.10 cm,
        partopsep=0pt,
        itemsep=0pt,
        leftmargin=10pt
    ]
}{
    \end{itemize}
}

\newenvironment{onecolentry}{
    \begin{adjustwidth}{0 cm + 0.00001 cm}{0 cm + 0.00001 cm}
}{
    \end{adjustwidth}
}

\newenvironment{twocolentry}[2][]{
    \onecolentry
    \def\secondColumn{#2}
    \setcolumnwidth{\fill, 4.5 cm}
    \begin{paracol}{2}
}{
    \switchcolumn \raggedleft \secondColumn
    \end{paracol}
    \endonecolentry
}

\newenvironment{threecolentry}[3][]{
    \onecolentry
    \def\thirdColumn{#3}
    \setcolumnwidth{, \fill, 4.5 cm}
    \begin{paracol}{3}
    {\raggedright #2} \switchcolumn
}{
    \switchcolumn \raggedleft \thirdColumn
    \end{paracol}
    \endonecolentry
}

\newenvironment{header}{
    \setlength{\topsep}{0pt}\par\kern\topsep\centering\linespread{1.5}
}{
    \par\kern\topsep
}

% "Last updated" watermark (optional; not printed unless you place it)
\newcommand{\placelastupdatedtext}{%
  \AddToShipoutPictureFG*{%
    \put(
        \LenToUnit{\paperwidth-2 cm-0 cm+0.05cm},
        \LenToUnit{\paperheight-1.0 cm}
    ){\vtop{{\null}\makebox[0pt][c]{%
        \small\color{gray}\textit{Последнее обновление: октябрь 2025}\hspace{\widthof{Последнее обновление: октябрь 2025}}%
    }}}%
  }%
}

% save the original href command in a new command:
\let\hrefWithoutArrow\href

\begin{document}
    \newcommand{\AND}{\unskip
        \cleaders\copy\ANDbox\hskip\wd\ANDbox
        \ignorespaces
    }
    \newsavebox\ANDbox
    \sbox\ANDbox{$|$}

    % ===== Header =====
    \begin{header}
        \fontsize{25 pt}{25 pt}\selectfont Shaxriyor Jabborov

        \vspace{5 pt}

        \normalsize
        \mbox{Ташкент}%
        \kern 5.0 pt%
        \AND%
        \kern 5.0 pt%
        \mbox{\hrefWithoutArrow{mailto:ilyshaxa@gmail.com}{ilyshaxa@gmail.com}}%
        \kern 5.0 pt%
        \AND%
        \kern 5.0 pt%
        \mbox{\hrefWithoutArrow{https://shaxa.dev}{shaxa.dev}}%
        \kern 5.0 pt%
        \AND%
        \kern 5.0 pt%
        \mbox{\hrefWithoutArrow{tel:+998-93-766-50-10}{+998 93 766 50 10}}%
        \kern 5.0 pt%
        \AND%
        \kern 5.0 pt%
        \mbox{\hrefWithoutArrow{https://linkedin.com/in/shaxriyor}{linkedin.com/in/shaxriyor}}%
        \kern 5.0 pt%
        \AND%
        \kern 5.0 pt%
        \mbox{\hrefWithoutArrow{https://github.com/ilyshaxa}{github.com/ilyshaxa}}%
    \end{header}

    \section{О себе}
    \begin{onecolentry}
        Cloud и DevOps инженер, специализирующийся на безопасных распределенных системах, оптимизации облачной инфраструктуры и автоматизации CI/CD промышленного уровня. Имеет опыт работы в средах AWS и GCP, в инженерии надежности на основе наблюдаемости, а также в безопасных гибридных интеграциях устройств с облаком. Подтвержденный опыт снижения инфраструктурных расходов, ускорения деплоев и усиления безопасности корпоративных систем.
    \end{onecolentry}

    \vspace{5 pt - 0.3 cm}

    % ===== Experience =====
    \section{Опыт работы}

    % KPI.com
    \begin{twocolentry}{May 2023 -- Oct 2026}
        \textbf{DevOps Engineer}, kpi.com -- Tashkent, UZ
    \end{twocolentry}
    \vspace{0.10 cm}
    \begin{onecolentry}
        \begin{highlights}
            \item Спроектировал и внедрил централизованный стек наблюдаемости (Prometheus, Grafana, Loki, Alertmanager), повысив прозрачность системы и сократив время реакции на инциденты на 60\%
            \item Разработал и реализовал безопасную облачную интеграцию системы учета посещаемости с использованием устройств распознавания лиц Hikvision по ISUP 5.0 (исходящий TLS), устранив входящую экспозицию на площадках клиентов и обеспечив безопасные API-интеграции
            \item Создал кастомный экспортер Prometheus на Python для публикации бизнес-KPI в режиме реального времени и принятия решений на основе данных
            \item Автоматизировал синхронизацию схем баз данных между dev/stage/prod через CLI-инструмент на Python, снизив количество ручных ошибок и операционные затраты
            \item Оптимизировал AWS-инфраструктуру (EC2, S3) и внедрил Savings Plans, сократив ежемесячные облачные расходы на 20\%
            \item Предотвратил сбои в продакшене, выявив узкие места из-за исчерпания потоков Tomcat и внедрив механизмы автоматического восстановления
            \item Усилил безопасность и производительность CI/CD, обновив устаревший GitLab Omnibus до последнего стабильного релиза
        \end{highlights}
    \end{onecolentry}
    \vspace{0.20 cm}
    
    \begin{twocolentry}{Jan 2025 -- July 2026}
        \textbf{DevOps Engineer}, PraaktisGo -- Tashkent, UZ
    \end{twocolentry}
    \vspace{0.10 cm}
    \begin{onecolentry}
        \begin{highlights}
        \item Руководил миграцией production AWS-инфраструктуры в новый AWS-аккаунт и регион без простоя, обеспечив полную целостность данных и непрерывную доступность сервисов
        \item Достиг 99.9\% доступности системы за счет внедрения проактивного мониторинга, алертинга и механизмов предотвращения инцидентов
        \item Сократил время деплоя на 50\% благодаря автоматизации CI/CD с Jenkins и стандартизации окружений с помощью Docker
        \item Спроектировал и внедрил стратегию кросс-облачного учета затрат и мониторинга производительности для долгосрочной масштабируемости гибридного облака
        \item Привлек \$100K кредитов Google Cloud (программа GFS Cloud Program 2024) и \$10K кредитов AWS Activate Portfolio для ускорения расширения инфраструктуры и экспериментов
        \item Оптимизировал инфраструктурные расходы и снизил операционные издержки на 15\% за счет стратегического использования Google Workspace Business Plan и практик управления облачными затратами
        \item Внедрил лучшие практики безопасности в средах AWS и GCP, включая усиление IAM, сегментацию сети и политики контролируемого доступа
        \item Задокументировал облачную архитектуру, процессы деплоя и операционные стандарты для ускорения онбординга команды и обеспечения непрерывности систем
\end{highlights}
\end{onecolentry}
    \vspace{0.20 cm}
    
    % DataSite Technology
    \begin{twocolentry}{Sep 2022 -- May 2023}
        \textbf{Frontend Engineer (Intern)}, DataSite Technology -- Tashkent, UZ
    \end{twocolentry}
    \vspace{0.10 cm}
    \begin{onecolentry}
        \begin{highlights}
            \item Разрабатывал адаптивные веб-интерфейсы с использованием HTML5, CSS3, JavaScript и React
            \item Преобразовывал дизайн-макеты Figma в функциональные страницы, готовые к production
            \item Сотрудничал с командами backend и DevOps, получив ранний практический опыт в процессах деплоя и практиках контроля версий
        \end{highlights}
    \end{onecolentry}
    \vspace{0.90 cm}
    
    % ===== Education =====
    \section{Образование}
    \begin{twocolentry}{2022 -- 2027}
        \textbf{Tashkent State University of Economics}, Bachelor's Degree in \textbf{Economics}
    \end{twocolentry}
    
\section{Технологии}

\begin{onecolentry}
\textbf{Cloud:} AWS, GCP
\end{onecolentry}

\vspace{0.2 cm}

\begin{onecolentry}
\textbf{Контейнеры и оркестрация:} Docker
\end{onecolentry}

\vspace{0.2 cm}

\begin{onecolentry}
\textbf{CI/CD:} Jenkins, GitLab CI, GitHub Actions
\end{onecolentry}

\vspace{0.2 cm}

\begin{onecolentry}
\textbf{Наблюдаемость:} Prometheus, Grafana, Loki, Alertmanager, ELK, Datadog
\end{onecolentry}

\vspace{0.2 cm}

\begin{onecolentry}
\textbf{Инфраструктура и безопасность:} Linux, Nginx, Apache Tomcat, Vault, Networking, Load Balancing, DNS
\end{onecolentry}

\vspace{0.2 cm}

\begin{onecolentry}
\textbf{Базы данных:} PostgreSQL, MySQL, MongoDB, Redis
\end{onecolentry}

\vspace{0.2 cm}

\begin{onecolentry}
\textbf{Языки программирования:} Python, Bash
\end{onecolentry}

    % ===== Languages =====
    \section{Языки}
    \begin{onecolentry}
        Узбекский (родной) \ \AND\ Английский (C1 - Продвинутый) \ \AND\ Русский (A1 - Базовый)
    \end{onecolentry}

\end{document}
